\subsection{A variable error bound for the good projector}

Let $G=G_0\times G_1$ be the residue group at the two level primes. Regard
each factor as a subgroup of $G$. Write $X$ for the complex Hilbert space
with orthonormal basis indexed by the fitted classes and Bessel coordinates,
$d_X=\dim X>0$, and $Y$ for the principal-level target with its Petersson
metric. Write $C$ for the isomorphism from source forms to these coordinates.
If $J$ is the principal level, $\nu>0$ the fitted unipotent covolume and
$\Pi$ the Poincar\'e map, fix the single operator
\[
       P=(\operatorname{vol}(J)\nu)^{-1/2}\Pi C^{-1},\qquad
       q_g=P^*\rho(g)P ,
\]
where $\rho$ is the unitary residue-group action on $Y$. Both moment
inequalities use this $P$. The unfolding needed here identifies each block
of $q_g$ with the conjugate of its full amplitude sum divided by $\nu$.
At the level of scalars this follows from the fitted Petersson identity
with factor $\operatorname{vol}(J)$ for the unscaled map, since the inner
product uses two copies of the scaling factor. The period-to-series identity
and the Petersson identity for the actual fitted functions are hypotheses
of this identification. The source gives the
\lineSrc{New/Spanning/SecondMoment/SpanningTheorem}{153}{coordinate map and scaling} and
the \lineSrc{New/Spanning/SecondMoment/SpanningTheorem}{168}{unfolded block identity}.

Let $R\subseteq G$ be the reduced Bessel subgroup and $r$ its unitary
representation on $X$, extended by zero off $R$. Assume the fitted
equivariance $\rho(s)P=Pr(s)$ for $s\in R$. Set $E_g=q_g-r(g)$.
For $s\in R$ we then have
$E_{sg}=r(s)E_g$ and $E_{gs}=E_g r(s)$. Define nonnegative errors by
\[
 h_R^2=\frac{1}{|R|d_X}\sum_{s\in R}\|E_s\|_{\mathrm{HS}}^2,\qquad
 h_i^2=\frac{1}{|R|d_X}\sum_{g\in R G_i}\|E_g\|_{\mathrm{HS}}^2
 \quad (i=0,1).
\]
For $t\in G_i$, assume the diagonal Bessel trace identity
$\operatorname{tr}r(t)=d_X\psi_i(t)$, with $\psi_i$ the reduced
unipotent character extended by zero outside its support. Write
\[
  \overline\psi_i^{\,\mathrm{cl}}(t)
     = |G_i|^{-1}\sum_{v\in G_i}\psi_i(vtv^{-1}),\qquad
  K_i=\sum_{t\in G_i}|\overline\psi_i^{\,\mathrm{cl}}(t)|^2 .
\]
The character average is normalized; the sum defining $K_i$ is unnormalized.
Conjugating $t$ by $G$ gives the same average because the other factor
commutes with $t$. Hence the full-group bound (10.11) implies $K_i\le2$.

Write $e_i$ for the self-adjoint good central idempotent in $\mathbb C[G_i]$ and
$x_i=1-e_i=\sum_{t\in G_i}x_i(t)t$. Put $a_i=x_i(1)$ and
$Q=\rho(e_0)\rho(e_1)$. The coefficients of $e_0\otimes e_1$ at
$(g_0,g_1)$ are the products of the two sums over allowed label tuples
having a primitive member at their respective primes. Expansion sums over
precisely the good label families appearing in the actual target projector.
This proves equality of the two projectors, coefficient by coefficient, for
the same label idempotents. See the
\lineSrc{New/Spanning/SecondMoment/SpanningTheorem}{197}{tensor identity}.

\begin{proposition}\label{prop:variable-projector-loss}
With the identities and conventions above, every subspace $Z\subseteq X$
satisfies
\[
 \|QP\operatorname{pr}_Z\|_{\mathrm{HS}}^2
 \ \ge\ \dim Z-d_X\left(h_R+
     \sum_{i=0}^1\bigl(h_i+\sqrt{a_i K_i}\bigr)\right).
\]
\end{proposition}
\begin{proof}
For $s\in R$, $E_s=(P^*P-1)r(s)$, so
$\|P^*P-1\|_{\mathrm{HS}}=h_R\sqrt{d_X}$. Pairing with
$\operatorname{pr}_Z$ in the Hilbert--Schmidt inner product gives
$\|P\operatorname{pr}_Z\|_{\mathrm{HS}}^2\ge\dim Z-h_R d_X$.

Fix $i$ and put $x=x_i$. On $\ell^2(G;\operatorname{End}(X))$, with
counting measure and Hilbert--Schmidt fibres, set
$(T_x F)(g)=\sum_{t\in G_i}x(t)F(gt)$. The right translations are unitary
and form a representation, so $T_x$ is an orthogonal projection.
Apply it to $F(g)=\mathbf1_{R G_i}(g)E_g$. For each $s\in R$,
$(T_xF)(s)=r(s)\sum_t x(t)E_t$. Consequently
\[
  |R|\Big\|\sum_t x(t)E_t\Big\|_{\mathrm{HS}}^2
     \le \|T_x F\|_2^2\le\|F\|_2^2=|R|d_Xh_i^2 .
\]
The absolute trace of $\sum_t x(t)E_t$ is
therefore at most $d_X h_i$. For the Bessel part, centrality of $x$ permits
averaging $\psi_i$ over conjugacy classes. The identity coefficient of
$x^*x=x$ gives $\sum_t|x(t)|^2=a_i$. Thus
\[
  \left|\sum_t x(t)\operatorname{tr}r(t)\right|
        \le d_X\sqrt{a_i K_i}.
\]
Adding the two contributions bounds
$\|\rho(x_i)P\|_{\mathrm{HS}}^2
 =\operatorname{tr}(P^*\rho(x_i)P)$ by $d_X(h_i+\sqrt{a_iK_i})$.
Since the two projections commute, $1-Q\le\rho(x_0)+\rho(x_1)$ as
positive operators. Restricting the source to $Z$ can only decrease the
squared Hilbert--Schmidt norm of the discarded part. Subtracting the two
losses from the first bound proves the assertion.
\end{proof}

For example, the arithmetic bounds $h_i\le b^{-6}$, $a_i\le b^{-8}/2$
and $K_i\le2$ imply (10.2), since $h_R\le h_i$ and
$3b^{-6}+2b^{-4}\le5b^{-4}$ for $b\ge1$.
On the constructed scales, if the former bounds for $h_i,a_i$ hold, it
suffices instead that $K_i=O(b^{6-\delta})$ for a fixed $\delta>0$,
with uniform implied constants at the fixed parameters. Indeed the total
loss is $O(b^{-6}+b^{-1-\delta/2})=o((32b x^{202})^{-1})$ because
$\log x=o(\log b)$. This relaxes the character input specifically for
the projected inequality. The full moment still has its own hypotheses.

For comparison, the fitted source invokes the
\lineSrc{New/Spanning/Unfolding/Moments}{446}{off-Bessel square sum},
the \lineSrc{New/Spanning/Unfolding/Moments}{751}{diagonal trace}, the
\lineSrc{New/Spanning/Unfolding/Moments}{991}{character mass} and the
\lineSrc{New/Spanning/SecondMoment/SpanningTheorem}{221}{bad idempotent mass}.
They give the
\lineSrc{New/Spanning/SecondMoment/SpanningTheorem}{257}{projected estimate}
for $QP$. The
\lineSrc{New/Spanning/SecondMoment/SpanningTheorem}{270}{full estimate}
uses the same $q_g$ and relabels the product group by its chosen places,
a bijection preserving the unnormalized sum.
